\vspace{-3mm}
\section{\textls[0]{A Holistic View of the Hallucination Spectrum: its Types and Scales}}
\label{sec:cat_hallucination}
\vspace{-2mm}
The issue of hallucination garnered research attention as early as \cite{maynez-etal-2020-faithfulness}. However, with the growing size of LLMs (\emph{empirical evidence provided in \cref{sec:mitigation}}), there is a corresponding increase in LLMs' susceptibility to hallucination. Consequently, there is a growing interest within the research community to study and understand hallucination to design mitigation techniques.

Researchers have loosely defined hallucinations and studied various notions of hallucinations in isolation.  Early exploration of factual vs. non-factual prompts for checking factuality of LMs is addressed in \cite{NEURIPS2022_df438caa}. A recent survey conducted by \citep{maynez-etal-2020-faithfulness}, categorized hallucination into two limited classes: intrinsic and extrinsic. Another recent paper \citep{ladhak2023pre}, delved into an intriguing type of hallucination known as name-nationality %\ad{one sentence to define what is it?} 
category hallucination. Several other papers \cite{raunak-etal-2021-curious,maynez-etal-2020-faithfulness} have explored task-specific categories of hallucination, such as summarization, question-answering, machine translation, etc.  Preliminary exploration of factual versus non-factual prompts is also addressed in \cite{NEURIPS2022_df438caa}. However, we purposefully avoid confining our study to a specific task to study the nuances of hallucination. Our main contention is that hallucination can occur in any NLG task, necessitating a thorough examination based on the fundamental principles of text generation from a given prompt. The findings from this study can be applied and extended to various NLP tasks. Thus, this paper aims to offer a comprehensive categorization of hallucination, as outlined below (see \cref{fig:types}).
%Some of the previous works have tried to categorize hallucination into name-nationality category \citep{ladhak2023pre}. O Although such broad categories of hallucination have been proposed, there is a lack of a finer-grained comprehensive set of categories that can help delineate the nuances. Put simply, while existing literature has defined hallucination in bits and pieces, we believe we are the first to provide a holistic treatment of the topic -- spanning the orientation, degree, and type of hallucination -- along with a large-scale dataset to further research in this area. \textbf{we are defining hallucination more comprehensive way:} Even if the above-mentioned works have defined hallucination patterns, in this work, we are defining hallucination in a more comprehensive way.
%The emerging capabilities of LLMs are receiving accolades; however, the associated undesired by-product is hallucination. In other words \textit{the emerging abilities of eloquent falsification}. 

%The NLP-AI community has recognized the significance \cite{} of comprehending unwanted hallucinations and has emphasized the need for developing improved techniques to detect and mitigate such occurrences in LLMs, in order to enhance their trustworthiness and usability. Although the previously mentioned research endeavors have explored and documented certain aspects of hallucination, this paper presents a thorough classification of hallucination generated by LLMs. Our scientific inquisitiveness was fueled by the following set of queries:

%Here we define a comprehensive spectrum of LLM-generated hallucinations:
\vspace{-3mm}
\subsection{Orientations of Hallucination}
We introduce two primary orientations of hallucination: (i) Factual Mirage (FM) and (ii) Silver Lining (SL), defined and exemplified below.
%Subsequent sections delve into detailed definitions and offer illustrative examples for each of these orientations and their subcategories.

\vspace{-1mm}
\subsubsection{Factual Mirage} 
Factual mirage (FM) is defined as the phenomenon wherein an LLM engages in hallucination or distortion of a  given prompt that is \ul{factually correct}. FM can be subdivided into two distinct sub-categories.

\vspace{-2mm}
\begin{tcolorbox}[enhanced,attach boxed title to top right={yshift=-2.5mm,yshifttext=-1mm},
  left=1pt,right=1pt,top=1pt,bottom=1pt,colback=blue!5!white,colframe=blue!75!black,colbacktitle=cyan!80!black,
  title=mild,fonttitle=\ttfamily\bfseries\scshape\fontsize{9}{9.6}\selectfont,
  boxed title style={size=fbox,colframe=cyan!50!black} ]
  %\vspace{-2mm}
  \begin{spacing}{0.75}
  \textbf{\scriptsize Prompt:} \textit{\scriptsize Capital of France}
  \end{spacing}

  \vspace{-2mm}
  \DrawLine
  
  \begin{spacing}{0.75}
  \textbf{\scriptsize AI-generated text:} {\fontfamily{lmss}\scriptsize\selectfont
  ...Paris is also the world fashion capital...}
  \end{spacing}

  \vspace{-2mm}
  \DrawLine

  %\begin{spacing}{0.75}
  \textbf{\scriptsize Fact:} \scriptsize Paris.
  %\end{spacing}

  % \vspace{-1mm}

\end{tcolorbox}
\vspace{-4mm}

% \textbf{Innate factual mirage (IFM):}
\subparagraph{Intrinsic factual mirage (IFM):}
In the following example, the LLM is providing a correct response while adding additional supplementary facts such as ``\emph{the world fashion capital},'' resulting in distortion or hallucination, has also been described in \cite{cao2022hallucinated}. 
%For example:
%Consider the prompt ``\emph{What is the capital of France?}'' The expected AI-generated response should be ``\emph{Paris},'' but let's say the generated response is ``\emph{Paris is also the world fashion capital.}`` In this case, the AI is providing a correct response while adding additional supplementary facts such as ``the world fashion capital,'' resulting in distortion or hallucination, as described in \cite{cao2022hallucinated}.



\subparagraph{Extrinsic factual mirage (EFM):} EFM refers to the phenomenon where an LLM deviates from factual accuracy. For example: 
%This can be observed when the generated text contains information or statements that are not supported by facts. 
%For instance:
%, if we consider the prompt ``\emph{What is the capital of France?}'' and the AI-generated response is ``\emph{Tokyo},'' it exemplifies an instance of EFM.

\vspace{-3mm}
\begin{tcolorbox}[enhanced,attach boxed title to top right={yshift=-2.5mm,yshifttext=-1mm}, left=1pt,right=1pt,top=1pt,bottom=1pt,
  colback=blue!5!white,colframe=blue!75!black,colbacktitle=red!80!black,
  title=alarming,fonttitle=\ttfamily\bfseries\scshape\fontsize{9}{9.6}\selectfont,
  boxed title style={size=fbox,colframe=red!50!black} ]
    
    %\vspace{-2mm}
  \begin{spacing}{0.75}
  \textbf{\scriptsize Prompt:} \textit{\scriptsize Engineering effort to build Eiffel tower}
  \end{spacing}

  \vspace{-2mm}
  \DrawLine

  \begin{spacing}{0.75}
  \textbf{\scriptsize AI-generated text:} {\fontfamily{lmss}\scriptsize\selectfont
  ...Designed by Gustave Eiffel, it was inaugurated in 1889 to celebrate the 100\textsuperscript{th} anniversary of the European Civil War...}
  \end{spacing}

  \vspace{-2mm}
  \DrawLine

  %\begin{spacing}{0.75}
  \textbf{\scriptsize Fact:} \scriptsize Eiffel tower was built to celebrate the 100th anniversary of the French Revolution.
  %\end{spacing}

  %\vspace{-1mm}
\end{tcolorbox}
%\vspace{-1mm}

%\noindent
\subsubsection{Silver Lining (SL)} Silver lining (SL)
% referred to as Illustrious Fabrication, 
is defined as the phenomenon in which an LLM indulges in hallucination by conjuring an elaborate and captivating narrative based on a given prompt that is \ul{factually incorrect}. 
%While Garbage In Garbage Out (GIGO) is a common phenomenon with technology -- where an incorrect input could lead to an unexpected outcome -- however, this is particularly worrisome in the context of LLMs due to the lack of a signal that indicates that the output is not factually consistent and could contain figments of the LLM's imagination. For example -
%given a prompt such as \emph{``Kamala Harris and Elon Musk are getting married,''} the LLM fabricates a story in response as follows: \emph{``Elon Musk got married to Kamala Harris in California on Thursday. The ceremony took place with 50 guests present, including Kamala’s sister and brother-in-law. Apart from Joe Biden all were present from the White House. Kamala’s ex-husband played the saxophone at the ceremony. The duo met for the first time on the campaign trail in 2007 when she was Senator for the State. Elon proposed to Kamala Harris in August 2019 on a beach. Kamala accepted and has been dating him ever since.''}.
%\acil{Why is this a ``silver lining'' if the generation (and prompt) is untrue? If we're not sure why, need to contemplate rephrasing the name.}
% \textbf{\acb{Internal} gilded phantasm (ISL)}
\vspace{-1mm}
\subparagraph{Intrinsic silver lining (ISL):} \hspace{0pt} ISL is the category when in some cases LLM does not generate a convincing story. As an example: %\ad{dont like the example!}
%For example: \ad{not happy about this example - not concrete}
%given a prompt \emph{Vladimir Putin has been elected president of the United States} if an LLM generates ``\emph{Vladimir Putin has been elected president of the United States. The Russian president-elect pledged to restore order and prosperity to America, following years of political unrest and economic decline}''.
\vspace{-4mm}
\begin{tcolorbox}[enhanced,attach boxed title to top right={yshift=-2.5mm,yshifttext=-1mm},
left=1pt,right=1pt,top=1pt,bottom=1pt,  colback=blue!5!white,colframe=blue!75!black,colbacktitle=violet!80!black,
  title=moderate,fonttitle=\ttfamily\bfseries\scshape\fontsize{9}{9.6}\selectfont,
  boxed title style={size=small,colframe=violet!50!black} ]

  \begin{spacing}{0.75}
  \textbf{\scriptsize Prompt:} \textit{\scriptsize Cell phone numbers go public this month. }
  \end{spacing}

  \vspace{-2mm}
  \DrawLine

  \begin{spacing}{0.75}
  \textbf{\scriptsize AI-generated text:} {\fontfamily{lmss}\scriptsize\selectfont
 ...Starting this month, cell phone numbers will be made public as a result of the implementation of the Right to Information (RTI) policy...}
  \end{spacing}

  \vspace{-2mm}
  \DrawLine

  %\begin{spacing}{0.75}
  \textbf{\scriptsize Fact:} \scriptsize This statement is false.
  %\end{spacing}

  %\vspace{-1mm}
  
\end{tcolorbox}
\vspace{-4mm}

\subparagraph{Extrinsic silver lining (ESL):} When an LLM generates a highly detailed and persuasive narrative in response to a factually incorrect prompt, it falls under the category of Extrinsic Silver Lining. % (ESL).  
%As an example:
% Please refer to the \adt{previous example} of Kamala Haris and Elon Musk prompt.


\vspace{-3mm}
\begin{tcolorbox}[enhanced,attach boxed title to top right={yshift=-2.5mm,yshifttext=-1mm},
left=1pt,right=1pt,top=1pt,bottom=1pt,  colback=blue!5!white,colframe=blue!75!black,colbacktitle=red!80!black,
  title=alarming,fonttitle=\ttfamily\bfseries\scshape\fontsize{9}{9.6}\selectfont,
  boxed title style={size=small,colframe=red!50!black} ]

  \begin{spacing}{0.75}
  \textbf{\scriptsize Prompt:} \textit{\scriptsize Kamala Harris and Elon Musk are getting married}
  \end{spacing}

  \vspace{-2mm}
  \DrawLine
  
  \begin{spacing}{0.75}
  \textbf{\scriptsize AI-generated text:} {\fontfamily{lmss}\scriptsize\selectfont
  Elon Musk got married to Kamala Harris in California on Thursday. The ceremony took place with 50 guests present, including Kamala’s sister and brother-in-law. Apart from Joe Biden all were present from the White House. Kamala’s ex-husband played the saxophone at the ceremony. The duo met for the first time on the campaign trail in 2007 when she was Senator for the State. Elon proposed to Kamala Harris in August 2019 on a beach. Kamala accepted and has been dating him ever since...}
  \end{spacing}
  

  \vspace{-2mm}
  \DrawLine
  
  %\begin{spacing}{0.75}
  \textbf{\scriptsize Fact:} \scriptsize The incident never happened!
  %\end{spacing}

  %\vspace{-2mm}
\end{tcolorbox}

%\vspace{-2mm}
\subsection{Categories of Hallucination}
\vspace{-1mm}
We meticulously classify hallucination into six distinct types, defined and exemplified below.

\subparagraph{Numeric Nuisance (NN):} This issue occurs when an LLM generates numeric values related to past events, such as dates, ages, or monetary amounts, that are inconsistent with the actual facts. %For instance:
% \textcolor{red}{Swagata: I feel like this is an example of time wrap trouble where the date of birth is wrong, this may not be an appropriate example of numeric nuisance.}

\vspace{-3mm}
\begin{tcolorbox}[enhanced,attach boxed title to top right={yshift=-2.5mm,yshifttext=-1mm},
left=1pt,right=1pt,top=1pt,bottom=1pt,  colback=blue!5!white,colframe=blue!75!black,colbacktitle=violet!80!black,
  title=moderate,fonttitle=\ttfamily\bfseries\scshape\fontsize{9}{9.6}\selectfont,
  boxed title style={size=small,colframe=violet!50!black} ]
%\vspace{-1mm}
  \begin{spacing}{0.75}
  \textbf{\scriptsize Prompt:} \textit{\scriptsize 
  Twitter is said to have laid off more workers.}
 \end{spacing}

  \vspace{-2mm}
  \DrawLine

  \begin{spacing}{0.75}
  \textbf{\scriptsize  AI-generated text:} {\fontfamily{lmss}\scriptsize\selectfont
  The layoffs come after Twitter announced earlier this month that it would be cutting its global workforce by 8\% or roughly 600 people...}
  \end{spacing}

  \vspace{-2mm}
  \DrawLine

  %\begin{spacing}{0.75}
  \textbf{\scriptsize  Fact:} \scriptsize Twitter laid off 200 people which is roughly 10\%
  %\end{spacing}

  % \vspace{-1mm}
\end{tcolorbox}
\vspace{-4mm}

\subparagraph{Acronym Ambiguity (AA):} This issue pertains to instances in which LLMs generate an imprecise expansion for an acronym. 
%For example:
%, here \textit{``RLHF''} is expanded as \textit{``RLHF is Reward-free Learning from Human Feedback in reinforcement learning''}. 
%It would indicate that the model is conjecturing based on its understanding and generating inaccurate information.
%This is a problem where LLM generates wrong expansion for a given acronym. For example: \adt{example}.

%\ac{WWWW or above RLHF example?}

\vspace{-4mm}
\begin{tcolorbox}[enhanced,attach boxed title to top right={yshift=-2.5mm,yshifttext=-1mm},
left=1pt,right=1pt,top=1pt,bottom=1pt,  colback=blue!5!white,colframe=blue!75!black,colbacktitle=violet!80!black,
  title=moderate,fonttitle=\ttfamily\bfseries\scshape\fontsize{9}{9.6}\selectfont,
  boxed title style={size=small,colframe=violet!50!black} ]
%\vspace{-1mm}

\begin{spacing}{0.75}
  \textbf{\scriptsize
  Prompt:} 
  \textit{\scriptsize RLHF in reinforcement learning}
  \end{spacing}

  \vspace{-2mm}
  \DrawLine

  \begin{spacing}{0.75}
  \textbf{\scriptsize AI-generated text:} {\fontfamily{lmss}\scriptsize\selectfont
  ...RLHF is Reward-free Learning from Human Feedback in reinforcement learning...}
  \end{spacing}

  \vspace{-2mm}
  \DrawLine

  %\begin{spacing}{0.75}
  \textbf{\scriptsize Fact:} \scriptsize RLHF stands for \textit{ ``Reinforcement Learning from Human Feedback''}
  %\end{spacing}

  % \vspace{-1mm}
\end{tcolorbox}
\vspace{-4mm}

\subparagraph{Generated Golem (GG):} This issue arises when an LLM fabricates an imaginary personality in relation to a past event, without concrete evidence. 
%For example:

\vspace{-5mm}
\begin{tcolorbox}[enhanced,attach boxed title to top right={yshift=-2.5mm,yshifttext=-1mm},
left=1pt,right=1pt,top=1pt,bottom=1pt,  colback=blue!5!white,colframe=blue!75!black,colbacktitle=violet!80!black,
  title=moderate,fonttitle=\ttfamily\bfseries\scshape\fontsize{9}{9.6}\selectfont,
  boxed title style={size=small,colframe=violet!50!black} ]
%\vspace{-1mm}

  \begin{spacing}{0.75}
   \textbf{\scriptsize Prompt:} \textit{\scriptsize Earthquake in Turkey}
  \end{spacing}

  \vspace{-2mm}
  \DrawLine

  \begin{spacing}{0.75}
  \textbf{\scriptsize AI-generated text:} {\fontfamily{lmss}\scriptsize\selectfont
  ...One rescuer, Hasan Cetin, said he was motivated by the thought of the survivors he helped save. “The moment when a man or a child is pulled out of the ruins, the joy of seeing them come back to life is...}
  %,” Cetin told The Associated Press. “This is why we keep on going.”
  \end{spacing}

  \vspace{-2mm}
  \DrawLine

  %\begin{spacing}{0.75}
  \textbf{\scriptsize Fact:} \scriptsize There is no evidence that the person Hasan Cetin exists.
  %\end{spacing}

  % \vspace{-1mm}

\end{tcolorbox}
\vspace{-4mm}

\subparagraph{Virtual Voice (VV):} At times LLMs generate quotations attributed to either fictional or real characters without sufficient evidence to verify the authenticity of such statements. 
%For example: 

%\ac{@Vipula: can we add quotes by Mahatma Gandhi or Steve Jobs for the genuine character aspect and try to elicit the source of an imaginary quote for the fictitious quote aspect?}

\vspace{-5mm}
\begin{tcolorbox}[enhanced,attach boxed title to top right={yshift=-2.5mm,yshifttext=-1mm},
left=1pt,right=1pt,top=1pt,bottom=1pt,  colback=blue!5!white,colframe=blue!75!black,colbacktitle=red!80!black,
  title=alarming,fonttitle=\ttfamily\bfseries\scshape\fontsize{9}{9.6}\selectfont,
  boxed title style={colframe=red!50!black} ]
  
    %\begin{spacing}{0.75}
  \textbf{\scriptsize Prompt:} \textit{\scriptsize Pfizer Press Release on COVID-19 vaccine}
  %\end{spacing}

  \vspace{-3mm}
  \DrawLine

  \begin{spacing}{0.75}
  \textbf{\scriptsize AI-generated text:} {\fontfamily{lmss}\scriptsize\selectfont
  ...Pfizer emphasized that their vaccine demonstrated an impressive efficacy rate... Pfizer CEO said, ``This is a giant leap for humanity..''...}
  \end{spacing}

  \vspace{-2mm}
  \DrawLine

  %\begin{spacing}{0.75}
  \textbf{\scriptsize Fact:} \scriptsize Pfizer CEO never said this.
  %\end{spacing}

  % \vspace{-1mm}
  
\end{tcolorbox}
\vspace{-2mm}

In this instance, the expression \emph{giant leap for humanity} is quoted from Neil Armstrong's renowned historical statement upon stepping onto the moon. 
% The issue of referencing the source of hallucination has also been recently explored in the work of  \cite{agrawal2023language}.

\vspace{-2mm}
\subparagraph{Geographic Erratum (GE):} This problem occurs when LLMs generate an incorrect location associated with an event. For example here in the first of the following two examples, Pentagon is stated to be situated in India.
%For example:


%\acil{@Vipula: can we ask an LLM where a particular moment lies and get it to generate an incorrect answer?}    

\subparagraph{Time Wrap (TW):} This problem entails LLMs generating text that exhibits a mashed fusion of events from different timelines. 
%For instance, 
%when presented with a prompt like ``\textit{USA on Ukraine war},'' if an LLM produces the response \textit{``U.S. President Barack Obama says the U.S. will not put troops in Ukraine,''} it showcases a distinct case of hallucination. This is evident as 
in the following case (2\textsuperscript{nd} example) the U.S. president during the Ukraine-Russia war is Joe Biden, not Barack Obama, thus contradicting the factual reality.

\vspace{-2.5mm}
\begin{tcolorbox}[enhanced,attach boxed title to top right={yshift=-2.5mm,yshifttext=-1mm},
left=1pt,right=1pt,top=1pt,bottom=1pt,  colback=blue!5!white,colframe=blue!75!black,colbacktitle=red!80!black,
  title=alarming,fonttitle=\ttfamily\bfseries\scshape\fontsize{9}{9.6}\selectfont,
  boxed title style={size=small,colframe=red!50!black} ]
  
  \begin{spacing}{0.75}
  \textbf{\scriptsize Prompt:} \textit{\scriptsize Prime Minister of the United Kingdom visiting India}
  \end{spacing}

  \vspace{-2mm}
  \DrawLine

  \begin{spacing}{0.75}
  \textbf{\scriptsize AI-generated text:} {\fontfamily{lmss}\scriptsize\selectfont
  ...The visit of the United Kingdom's Prime Minister to India holds immense significance...During his journey, he will also pay a visit to the Pentagon where he is scheduled to meet with Smriti Irani, the Minister of...}
  \end{spacing}

  \vspace{-2mm}
  \DrawLine

  %\begin{spacing}{0.75}
  \textbf{\scriptsize Fact:} \scriptsize Pentagon is the defense headquarters of the USA, located in Washington DC, USA -- not in India.
  %\end{spacing}

  % \vspace{-1mm}
\end{tcolorbox}
%\vspace{-4mm}

\vspace{-4mm}
\begin{tcolorbox}[enhanced,attach boxed title to top right={yshift=-2.5mm,yshifttext=-1mm},
left=1pt,right=1pt,top=1pt,bottom=1pt, colback=blue!5!white,colframe=blue!75!black,colbacktitle=red!80!black,
  title=alarming,fonttitle=\ttfamily\bfseries\scshape\fontsize{9}{9.6}\selectfont,
  boxed title style={size=small,colframe=red!50!black} ]

  \vspace{-1.25mm}
  \begin{spacing}{0.75}
  \textbf{\scriptsize Prompt:} \textit{\scriptsize USA on Ukraine war}
  \end{spacing}

  \vspace{-2mm}
  \DrawLine

  \begin{spacing}{0.75}
  \textbf{\scriptsize AI-generated text:} 
  {\fontfamily{lmss}\scriptsize\selectfont ...U.S. President Barack Obama says the U.S. will not put troops in Ukraine...}
  \end{spacing}

  \vspace{-2mm}
  \DrawLine

  %\begin{spacing}{0.75}
  \textbf{\scriptsize Fact:} \scriptsize The actual U.S. president during the Ukraine-Russia war is Joe Biden. 
  %\end{spacing}

  % \vspace{-1mm}
\end{tcolorbox}
%\vspace{-2mm}


%\begin{comment}
\subsection{Degrees of Hallucination}

We annotate the degree of hallucination using three levels: \textit{mild, moderate,} and \textit{alarming} (labeled as 0, 1, and 2 respectively). \textit{Mild} indicates minor hallucination which is superficial in terms of its impact. \textit{Moderate} indicates a level of hallucination that introduces facts that are either fictitious or tangential to the topic at hand. \textit{Alarming} indicates added information pieces that bear a radical dissemblance from the topic fed via the prompt. Please refer to \cref{subsec:annotation} for more details.

% This is labeled as \textbf{2}.
%\end{comment}



\begin{comment}
With the recent and rapid advances in the areas of LLMs and Generative AI, the pre-eminent and ubiquitous concern is of \emph{hallucination}. We release large-scale first-of-its-kind human-annotated data with detailed annotations on - intrinsic vs. extrinsic hallucinations, and degree of hallucination, and we ask participants to come up with either black-box factuality Assessment and/or evidence-based fact-checking. First, we define categories of hallucinations: 





\textbf{Intrinsic Hallucination}: Intrinsic hallucination refers to the phenomenon when an LLM generates text that topically slightly deviates from the input and/or has a lack of grounding in reality.
%For instance, in the abstractive summarization task from Table 1\act{*which table 1?*}, the generated summary.
% “\textit{The first Ebola vaccine was approved in 2021}” contradicts the source content “\textit{The first vaccine for Ebola was approved by the FDA in 2019.}”
For example, given a prompt ''\textit{USA on Ukraine war}'' an LLM generates \textit{''U.S. President Barack Obama says the U.S. will not put troops in Ukraine''}. We can see a clear case of intrinsic hallucination as the US president during the Ukraine-Russia war is Joe Biden, not Barack Obama, contradicting the reality.

\textbf{Extrinsic Hallucination}: We define extrinsic hallucination to be the generated output from an LLM that cannot be verified, or contradicted from the source content provided as prompt. For example, we provide a prompt stating \textit{''North Korea has conducted six underground nuclear tests, and a seventh may be on the way.''}, and the resulting output generated by the model was \textit{''The international community has condemned North Korea’s nuclear tests, with the UN Security Council imposing a range of sanctions in response. The US and other world powers have urged North Korea to abandon its nuclear weapons program and to comply with international law.''} As the source input makes no reference to the U.S. or U.N. Security Council imposing any sanctions, the claimed imposition of sanctions cannot be verified from the input alone.
\textbf{Degree of Hallucination:}We define the \emph{degree} of hallucination by the Likert scale\footnote{\url{https://en.wikipedia.org/wiki/Likert_scale}}.

\textbf{Orientation of Hallucination:}Like/dislike? \faThumbsOUp  \faThumbsODown

\subsection{Positive Hallucination}
We define \emph{positive hallucination} when the model detects hallucination in the generated text. For instance, \emph{What is the capital of France?} and the AI-generated response is \emph{Tokyo}, hallucination is present.

%gpt
``Positive hallucination'' refers to a phenomenon where a language model generates text that contains information that is not present in the input or the training data. In other words, the model generates output that is not supported by the context or the knowledge it has learned from the training data. Positive hallucination can occur due to various reasons, such as overfitting, lack of relevant training data, or biases in the model's architecture or training process. Positive hallucination can be a problem in natural language processing, as it can lead to the generation of misleading or inaccurate text. However, it can also be beneficial in some cases, such as in text completion or generation tasks where the model is expected to generate novel and creative output.

\subsection{Negative Hallucination}
We define \emph{negative hallucination} when the model detects falsely detects hallucination in the generated text. For instance, \emph{What is the capital of France?} and the AI-generated response is \emph{Paris. Paris is also the world fashion capital.}. Here, it is generating a correct response and additionally some more useful facts \cite{cao2022hallucinated} like \emph{Paris is also the world fashion capital}. Hence, negative hallucination is present.

\subsubsection{Half true, half false, false}
The goal of the Truth-O-Meter is to reflect the relative accuracy of a statement. The meter has six ratings, in decreasing level of truthfulness\footnote{\url{https://tinyurl.com/ycxw9t9z}}

\noindent\textbf{TRUE:} The statement is accurate and there’s nothing significant missing. \\
\textbf{MOSTLY TRUE:} The statement is accurate but needs clarification or additional information.  \\
\textbf{HALF TRUE:} The statement is partially accurate but leaves out important details or takes things out of context.  \\
\textbf{MOSTLY FALSE:} The statement contains an element of truth but ignores critical facts that would give a different impression.  \\
\textbf{FALSE:} The statement is not accurate.  \\
\textbf{PANTS ON FIRE:} The statement is not accurate and makes a ridiculous claim.
\end{comment}